\begin{table}
\centering
\caption{Parameter sizes of InteractE and ComDensE.}
\label{table3}
\begin{tabular}{c|c|c}
\hline
\multicolumn{3}{c}{\textbf{The number of parameters}}\\
\hline
 \textbf{method} & \textbf{FB15k-237} & \textbf{WN18RR}\\
 \hline
 %ConvE & 5,092,127 & 10,306,329  \\
 % \hline
 InteractE & 18M & 60M \\
  \hline
 %ComDensE \textbf{512 $\times d$} & 129,037,479 & 56,952,841 \\
%\hline
 ComDensE \textbf{256 $\times d$} & 66M & 33M \\
\hline
 ComDensE \textbf{128 $\times d$} & 35M & 22M \\
\hline
\end{tabular}
\end{table}

\begin{table}
\centering
\caption{All other hyperparameters are same except row dimension of relation-aware and common dense matrices. $d$ denote column dimension of dense matrices. While ComDensE uses 256 $\times d$ matrices to both benchmark datasets, ComDensE using 128 $\times d$ matrices achieves a similar performance gain on FB15k-237.
}
\label{table4}
\begin{tabular}{c|c|c}
\hline
\multicolumn{3}{c}{\textbf{Various sizes of matrix in ComDensE}}\\
\hline
 & \textbf{FB15k-237} & \textbf{WN18RR}\\
 \hline
 \multicolumn{3}{c}{\textbf{256$\times d$}}\\
 \textbf{MRR} & .356 & .473 \\
 \textbf{HIT@10} & .536 & .538 \\
 \textbf{HIT@1} & .265 & .440\\
  \hline
 \multicolumn{3}{c}{\textbf{128$\times d$}}\\
 \textbf{MRR} & .354 & .453 \\
 \textbf{HIT@10} & .534 & .512 \\
 \textbf{HIT@1} & .263 & .423 \\
\hline
\end{tabular}
%}
\end{table} 

\begin{table*}[t]
\centering
\caption{Link prediction results of baseline models. \textbf{\underline{Bold with underline}} : the best performance, \textbf{bold} : the second best performance.}
\label{table5}
\begin{tabular}{cccc|ccc}
\hline
 & \multicolumn{3}{c}{\textbf{FB15k-237}} & \multicolumn{3}{c}{\textbf{WN18RR}}\\
 \hline
 \textbf{Methods} &
 \textbf{MRR} &
 \textbf{HIT@10} &
 \textbf{HIT@1} &
 \textbf{MRR} &
 \textbf{HIT@10} &
 \textbf{HIT@1} \\
 \hline
 TransE \cite{bordes2013translating} & .279 & .441  & - & .243 & .532 & -\\
 RotatE \cite{sun2019rotate} & .338 & .533 & .241 & \textbf{\underline{.476}} & \textbf{\underline{.571}} & .428\\
 \hline
 DistMult \cite{yang2014embedding} & .241 & .419 & .155 & .430 & .49 & .39\\
 ComplEx \cite{trouillon2016complex} & .247 & .428 & .158 & .44 & .51 & .41\\
 \hline
 RGCN \cite{schlichtkrull2018modeling} & .248 & .417 & .151 & - & - & -\\
 \hline
 ConvE \cite{dettmers2018convolutional} & .325 &  .501 & .237 & .43 & .52 & .40\\
 ConvTransE \cite{shang2019end} & .33 & .51 & .24 & .46 & .52 & .43\\
 SACN \cite{shang2019end} & .35 &  \textbf{\underline{.54}} & .26 & .47 & \textbf{.54} & .43\\
 InteractE \cite{vashishth2020interacte} & \textbf{.354} & .535 & \textbf{.263} & .463 & .528 & \textbf{.430}\\
 \hline
 ComDensE (proposed) & \textbf{\underline{.356}}  &  {\textbf{.536}} & \textbf{\underline{.265}} & \textbf{.473}  & .538 & \textbf{\underline{.440}}\\
 \hline
\end{tabular}
%}
\end{table*}

\begin{table*}
\centering
\caption{Link prediction results by relation type on FB15k-237 dataset for ConvE, InteractE, and ComDensE. Relations are categorized into one-to-one (1:1), one-to-many (1:N), many-to-one (N:1) and many-to-many (N:N). We observe that ComDensE is effective in capturing most relations compared to baselines.}
\label{table6}
\begin{tabular}{cc|ccc|ccc|ccc}
\hline
  \multicolumn{2}{c}{} & \multicolumn{3}{c}{\textbf{ConvE}} & \multicolumn{3}{c}{\textbf{InteractE}} &  \multicolumn{3}{c}{\textbf{ComDensE}}\\
 \hline
 \multicolumn{2}{c}{\textbf{Relation Type}} &
 \textbf{MRR} &
 \textbf{HIT@10} &
 \textbf{HIT@1} &
 \textbf{MRR} &
 \textbf{HIT@10} &
 \textbf{HIT@1} &
 \textbf{MRR} &
 \textbf{HIT@10} &
 \textbf{HIT@1} \\
 \hline
 & \textbf{1:1} & .374 &.505 & .068 & .386 & .547 & .245 & \textbf{.422}  & \textbf{.557} & \textbf{.349}\\
 Pred & \textbf{1:N} & .095 & .17 & .030 & \textbf{.106} & \textbf{.192} & \textbf{.043} & .084 & .181 & \textbf{.043}\\
 Head & \textbf{N:1} & .444 & .644 & .315 & \textbf{.466} & .647 & .369 & \textbf{.466} & \textbf{.649} & \textbf{.372}\\
 & \textbf{N:N} & .261 & .459 & .140 & .276 & \textbf{.476} & .164 & \textbf{.279} & \textbf{.476} & \textbf{.187}\\
 \hline
 & \textbf{1:1} & .366 & .51 & .052 & .368 & .547 & .229  & \textbf{.422} & \textbf{.563} & \textbf{.349}\\
 Pred & \textbf{1:N} & .762 & .878 & .658 & .777 & .708 & \textbf{.881} & \textbf{.779} & \textbf{.884} & .717\\
 Tail & \textbf{N:1} & .069 &.15 & .015 & .074 &.141 & .034 & \textbf{.084} & \textbf{.169} & \textbf{.043}\\
 & \textbf{N:N} & .375 & .603 & .237 & .395 & .617 & .272 & \textbf{.396} & \textbf{.618} & \textbf{.285}\\
 \hline
\end{tabular}
%}
\end{table*}

\section{Experiment}

\subsection{Datasets and Evaluation Settings}
ComDensE is evaluated on two KG benchmark datasets: FB15k-237 \cite{toutanova2015observed} and WN18RR \cite{dettmers2018convolutional}. These datasets are subsets of FB15k \cite{bordes2013translating} and WN18 \cite{bordes2013translating}, respectively. Specifically, inverse relations are removed from the original datasets so as to prevent direct inference via reversing triples. The statistics of benchmark datasets are summarized in Table \ref{table2}. 


We evaluate the performance of link prediction in the filtered setting. Specifically, we compute scores from all other candidate triples in the test triples not existing in the training, validation, or test set, where candidates are generated by corrupting subjects for objects for predicting object. We use Mean Reciprocal Rank (MRR) and Hits at N (HIT@N) which are standard evaluation metrics for these datasets and models are evaluated in our experiments. We train and evaluate 5 times and then average performance results for robust evaluation. We compare ComDensE to baseline models: TransE \cite{bordes2013translating}, RotatE \cite{sun2019rotate}, DistMult \cite{yang2014embedding}, ComplEx \cite{trouillon2016complex}, RGCN \cite{schlichtkrull2018modeling}, ConvE \cite{dettmers2018convolutional}, ConvTransE \cite{shang2019end}, SACN \cite{shang2019end} and InteractE \cite{vashishth2020interacte}


\subsection{Hyperparameter Setting}
We use Adam~\cite{kingma2014adam} optimizer and search the hyperparameters on the validation set. The ranges of the hyperparameters for the grid design space: (embedding dimension) $\in$ \{150, 200, 256, 300\}, (batch size) $\in$ \{128, 256\}, (learning rate) $\in$ \{0.001, 0.0001\}, (input dropout) $\in$ \{0.4, 0.5\}, (hidden dropout) $\in$ \{0.4, 0.5\}, (width $n \times $ dense matrix $W^h \in \Re^{nd_h \times d}$, $n) \in$ \{2, 5, 50, 100, 200\}.\\
\indent Finally, we use same hyperparameters as embedding dimension 256, batch size 128, learning rate 0.0001, input dropout 0.4, row dimension of matrix $d_h$ 256 and hidden dropout 0.5 for both FB15k-237 and WN18RR. However, since the width is a crucial element of designing ComDensE and it depends on dataset sensitively, we only search a different width, where our method contains width 2 $\times$ square matrix $W^h \in \Re^{2d_h \times d}$ for FB15k-237 and width 100 $\times$ square matrix $W^h \in \Re^{100d_h \times d}$ for WN18RR, where $d_e$ denotes the dimension of entity embedding, $d_r$ denotes the dimension of relation embedding and $d := d_e + d_r$.